\chapter{奶蛙前史：一场正在发生的失传媒体考据}\label{ch:04}

\begin{nlkeypoints}
  \begin{enumerate}[label=(\arabic*),leftmargin=2.4em]
    \item 奶蛙的起源目前是一条\keyw{未结案的证据链}：
          最系统的社群考据指向 2024 年 10 月 2 日星野 AI 平台
          “纲手【百豪】”所创建的“奶龙”智能体，
          但该智能体现已不可访问，考据材料自身
          声明“仅为推测、缺乏直接证据”。
    \item 起源考据经历三个可辨识阶段：
          \keyw{误指}（“季洪行说”，2024-03 起）、
          \keyw{线索}（星野 AI 水印与截图，2024-12 前后）、
          \keyw{发现与证伪并存}（2026-08 智能体确认 +
          复现失败）。三阶段构成“错误归因→证据积累→
          归因修正”的标准考据动力学。
    \item 本章同时示范两条纪律：
          其一，\keyw{不可复核性}必须被显式写入结论
          （“指向某来源”与“证明某来源”是两回事）；
          其二，\keyw{考据伦理}——本书不推断、
          不转述任何关于具体个人的身份、健康、
          生死与私德的内容。
  \end{enumerate}
\end{nlkeypoints}

第三章结束时留下了一个问题：
正典一侧的材料有 A 级来源支撑，闭环良好；
变种一侧，最好的材料是一份自陈“待改进、来源缺失”
的维基词条。本章不回避这个落差，反而把它当作
全书的方法论示范对象：
\keyw{当证据链未结案时，研究者应当如何诚实地写作}。
具体而言，本章做四件事：
把奶蛙起源考据拆成三个可辨识阶段；
讨论“不可复核性”这一认识论困境；
单独处理笑声音源这条第二证据链；
以及——这可能是本章最重要的一节——
说明为什么关于“创作者”的写作
必须有一条不能逾越的伦理线。

\section{什么是“考据”：方法与标准}\label{sec:04:method}

在进入材料之前，先给“考据”一个操作性定义，
因为它在本书中承担特殊的功能：

\begin{definition}[网络考据（本书用法）]\label{def:04:kaoju}
  网络考据指：由非职业研究者组成的社群，
  以公开可检索的网络材料（截图、帖子、智能体档案、
  时间戳、水印）为证据，对某一网络文化对象的
  来源、作者、演化路径进行追溯与归因的活动。
  其结论形态通常是\keyw{带置信度的指向}
  （“最可能的来源是 X，证据为……”），
  而非\keyw{证明}（“X 是来源，证毕”）。
\end{definition}

这个定义里最要紧的是结论形态的差异。
网络考据与法庭证据、档案学考据共享“证据—推理—结论”
的骨架，但它缺少三样东西：
(1) 完整的证据保全（原始载体随时消失）；
(2) 独立的质证环节（同一社群既是调查者又是受众）；
(3) 明确的“结案”标准（指向某来源后，
证据链是否闭合由谁判定？通常没有判定者）。
因此，本书引用考据材料时，一律使用
“考据指向……”“有考据认为……”的措辞，
并保留考据材料自身的限定语。
这是把定义~\ref{def:04:kaoju} 落进句式的直接结果。

\begin{nlmethod}[研究方法提示]
  本章的材料来源与可信度：
  核心材料来自失传媒体中文维基
  （\conf{C}，词条自陈“待改进”“来源缺失”）
  与一篇 2026 年 9 月的二手评述文章
  （\conf{C}）；第一版教材的叙述
  （\conf{C}）作为早期考据的快照保留。
  全书没有一条 \conf{A} 或 \conf{B} 级来源
  直接描述奶蛙起源——
  这一“缺口”本身就是本章要分析的现象之一
  （官方与主流媒体对民间二创的沉默）。
\end{nlmethod}

\section{第一阶段：误指（2024-03 起）}\label{sec:04:mistaken}

按失传媒体中文维基词条的整理，
最早可观测的“奶蛙”相关传播出现在 2024 年 3 月 7 日：
有网友将“奶蛙”形象与抖音博主“季洪行”的视频进行
AI 融合\conf{C}\cite{lostmedia2026naiwa}。
由于该融合内容在一段时间内传播最广，
“奶蛙”长期被认为是
“季洪行视频 + 奶龙 AI 融合”的产物
\conf{C}\cite{lostmedia2026naiwa}。

这个错误归因有三个值得注意的结构性原因：

\begin{enumerate}[label=\arabic*.]
  \item \textbf{可见性偏误}：
        被误指的对象是“当时最出圈”的内容。
        考据的对象如果是“第一个存在的”，
        调查者实际找到的往往只是
        “第一个被广泛看见的”——
        传播链的可见端点与起源端点并不重合。
  \item \textbf{时间标记的误导性}：
        融合内容出现时，“我是奶龙”梗
        （正典奶龙的“糖”内容，第一版教材有详细讨论）
        正在高峰期，奶龙形象的
        AI 加工需求本就旺盛，
        使得“融合说”在直觉上可信。
  \item \textbf{小范围传播的“真空期”}：
        词条的判断是：当时最出圈的是
        “我是奶龙”此梗而不是奶蛙，
        因此奶蛙“在小范围传播”，
        直到 2026 年“奶蛙捧腹大笑”
        视频才爆火\conf{C}\cite{lostmedia2026naiwa}。
        这一判断（若成立）意味着
        奶蛙在 2024—2025 年间
        存在一个“低可见度的潜伏期”——
        这是第五章传播模型里
        “二次爆发”假设（\ref{sec:05:sir}）
        的第一个经验支撑。
\end{enumerate}

\begin{nlcritique}[批判性反思]
  误指阶段还留下一个次生问题：
  当“错误答案”先于正确答案传播了
  一年半以上，它本身会变成一种
  \keyw{文化事实}——大量早期材料、
  表情包与口头叙述都以“季洪行融合说”为前提。
  考据修正“来源”时，
  无法回滚已经基于错误来源
  生产出来的二创。
  这在第七章的“形式稳定、内容可换”
  分析中会被反复用到。
\end{nlcritique}

\section{第二阶段：星野 AI 线索（2024-12 前后）}\label{sec:04:clue}

词条记录的第二阶段线索来自
抖音用户对账号“欢迎大家恐龙世”
（现名“欢迎大家恐龙世界”）的声讨：
该账号“经常发布儿童邪典内容”，
因用户声讨后被官方封禁、转私密账号
\conf{C}\cite{lostmedia2026naiwa}。
在声讨材料的截图中，
有人发现了“奶蛙”的照片
（声讨发布时间在 2024 年 12 月前后）
\conf{C}\cite{lostmedia2026naiwa}；
进一步地，该照片颜色反转后
可见 AI 生成水印，
结合配文“星野ai怎么了”，
用户推测其来源为星野 AI 生成
\conf{C}\cite{lostmedia2026naiwa}。

从证据学上看，这一段的价值在于
它把候选来源从“不确定的某个博主”
收窄到“某个具体平台的某个生成环节”——
\keyw{候选空间缩小}是考据推进的主要形态，
而不是“一步找到真凶”。
同时，这条线索也带有典型的
“孤证”特征：一张截图、一段配文，
均出自被声讨的账号生态，
且原始生成过程无法重放。
本书对此的处理是：
保留其作为“线索”的地位，
不赋予其“证据”的地位。

\section{第三阶段：发现与证伪并存（2026-08）}\label{sec:04:discovery}

2026 年 8 月 4 日，
B 站用户“雾蚀心忘症”发布评论，
表示在星野 App 中找到了
符合形象的奶蛙智能体
\conf{C}\cite{lostmedia2026naiwa}。
词条随后将智能体定位到
星野用户“纲手【百豪】”于
2024 年 10 月 2 日创建的
“奶龙”智能体，ID 为 \fname{zNCvQDKfXy}：
其背景图片（即头像）即最初版奶蛙形象；
“奶龙”智能体因 AI 渲染导致
体型畸变——腹部膨大、四肢细小、形似青蛙——
故得名“奶蛙”
\conf{C}\cite{lostmedia2026naiwa}。
该智能体“已经无法搜索及分享，
评论区也已关闭”
\conf{C}\cite{lostmedia2026naiwa}。

第一版教材记录的考据成果与之相互印证：
“2026 年 8 月 5 日，
被以 B 站 UP 主看仙碎为首的一大批
奶蛙爱好者解决（视频号 BV1ndMk65EwM）”
\conf{C}\cite{shoutanji2026naiwa}。
两个来源（维基词条与第一版教材/评述文章）
在核心事实上\keyw{一致}：
来源指向星野平台、
“纲手【百豪】”账号、
2024-10-02 的智能体、
2026-08 的确认时点。
一致性本身提高了指向的可信度——
但请注意，它提高的是
“两个 C 级来源指向同一处”的可信度，
而不是“该处即起源”的证明力。

\begin{nlhistory}[史料与考据]
发现之后，考据并未“结案”，而是进入了\keyw{证伪与复现}阶段。
 词条转述了 B 站用户“321BBnm”复现视频的结论：
 “‘奶蛙’生成来源仍存在诸多疑点，虽推测其诞生于星野 App，
 但该智能体已无法搜索，
 无法确认是星野 AI 直接生成还是其他 AI 生成后上传。
 多种生成路径的推测均被证伪，其多次复现也未成功，结合老用户所说：
 不管是现在还是原来，星野 AI 本身能力有限，
 无法生成出与参考图契合的图片。最终‘星野 App 出现 bug’的结论仅为推测，
缺乏直接证据，且相关智能体已无法访问，
 生成来源的信息完整性和可追溯性降低”\conf{C}\cite{lostmedia2026naiwa}。
\end{nlhistory}

这段引文是本章的核心教学材料。它示范了网络考据最诚实的写法：
 在给出指向的同时，逐条列出指向所依赖的每一个薄弱环节。
 本书在引用它时，连同这些限定一起引用——
 \keyw{不替考据者删掉他们自己写的“但是”}。

\section{不可复核性：一个认识论困境}\label{sec:04:irrevocable}

把第三阶段的引文推到底，会得到一个一般性的问题：当关键证据载体消失时，
 “来源指向”还能不能当作知识使用？

本书的立场分三层：

\begin{enumerate}[label=\arabic*.]
  \item \textbf{指向可以保留，证明力要降级}。
“考据指向星野平台某智能体”这一陈述，其真值依赖于一个已不可访问的对象；
 但支撑它的材料（截图、评论、时间戳）多数仍可公开检索。本书的做法：
 保留指向，把结论措辞从“奶蛙来自 X”降为“奶蛙最可能的来源指向 X，
 且该指向目前不可复核”。
  \item \textbf{不可复核性不是孤例，
而是网络文化的常态属性}。账号注销、智能体下架、视频删除、平台迁移——
 每一个环节都会让某段传播链的中间环节消失。
 第五章的传播模型把这一现象形式化为“证据损耗率”（M 级参数），
 以说明为什么网络文化研究必须以“现存样本”而非“原始总体”为对象。
  \item \textbf{对教学而言，未结案是优点}。
一个被完整证伪或完整证实的起源故事只能被转述；
 一个未结案的故事可以被用来训练证据分级、置信度措辞、
 复现与证伪的概念——这正是本章的存在理由。
\end{enumerate}

\begin{nltermbox}{术语：不可复核性（irrevocability，本书用法）}
指关键证据载体在指向形成\keyw{之后}消失，
 导致指向无法被独立第三方重新验证的状态。与“来源不明”（从未找到载体）
 相对：不可复核性意味着载体曾经被找到、被记录、随后消失。
 奶蛙智能体 \fname{zNCvQDKfXy} 是本书讨论的标准案例。
\end{nltermbox}

\section{第二证据链：捧腹大笑的音源}\label{sec:04:audio}

形象来源与声音来源是
两条独立的证据链，
本书严格分开处理。

按第一版教材的记录：主播“王麻衣”在直播时，一名自称“布洛芬”
的青年用户连麦，“连麦之后，青年不知为何一直发出异常爽朗的大笑声，
 把所有人都搞破绷了”；
 “如同黑雷霆的目力咆哮与野兽先辈一个一个一个表情的搭配一样，
 这经典的笑声便流传下来，成为了奶蛙捧腹大笑的灵魂音色”
\conf{C} （第一版教材原文，本书转引，级别维持 \conf{C}）。
 二手评述补充了传播端的描述：
 奶蛙二创视频“经常会与变调版不知火舞语音搭配”，
 二者合一后“在不同身份、不同职业的奶蛙们会一同捧腹大笑”
的固定结尾中反复出现 \conf{C}\cite{shoutanji2026naiwa}。

两条材料合起来，可以得到一个 \keyw{“声音—形象合流”} 的传播学描述：
 一个与形象无关的直播笑声素材，经过变调处理后，
 与奶蛙形象在视频二创的剪辑环节发生绑定；绑定之所以稳定，
 是因为“结尾集体大笑”
 这个形式本身具有独立的可传播性（第九章将把它当作“可传播性单位”分析）。

必须同时写明这条证据链的全部脆弱性：音源的具体直播场次、原始录音、
 “变调”发生的具体环节，目前均无可检索的直接材料；“王麻衣”“布洛芬”
 两个名称均出自单线来源（第一版教材）；本书对其的处理是：
 作为 \keyw{“社群追溯到的音源指向”} 记录，级别 \conf{C}，
 不作为任何经验结论的前提。第六章将给出“捧腹大笑”
作为 \keyw{声音符号}的独立分析——那部分分析不依赖音源考据是否成立。

\section{2026 年 3 月：关评、拉黑与“潘多拉奶盒”}\label{sec:04:official}

官方介入的时间点必须单独列出，因为它同时是产业线（第三章）
 与现象线的 \keyw{分界事件}。

第一版教材的表述：“2026 年 3 月，奶龙官方呼吁停止丑化二创，
 关闭视频评论区并批量拉黑发奶蛙的账号。
 但是奶蛙的影响力已经扩散开来（潘多拉奶盒打开了），难以制止了。
 至此，奶蛙凭借着魔性而意义不明的捧腹大笑，正式脱离了奶龙，
 成为了互联网上一个强而有力的热梗” \conf{C}（第一版教材原文，转引，
级别维持 \conf{C}）。

二手评述的交叉材料：“奶龙官方对于奶蛙梗已经深恶痛绝。
 之前网友一起玩梗还能被法不责众的幌子遮掩一下。
 而这时当奶蛙的源头有了明确的指向性时——
 只能说奶龙法务部的拳头想喝血了”（评述者语）
\conf{C} \cite{shoutanji2026naiwa}；同一评述同时记录：
 “直至今日你依然可以在奶龙官方视频的评论区看见无数奶蛙图片”
 \conf{C}\cite{shoutanji2026naiwa}。

本书从这三条材料中提取两个分析判断（M 级，非经验断言）：

\begin{enumerate}[label=\arabic*.]
  \item \textbf{抑制是“不彻底”的}：
关评、拉黑针对的是官方账号下的评论区与发布奶蛙的特定账号，
 而不是奶蛙形象本身在整个互联网上的流通——
 “官方视频评论区仍有无数奶蛙图片”
 说明抑制的执行在空间上就是不完整的。
  \item \textbf{“源头指向性”改变了
抑制的动机结构}：评述者把法务动作的时间点与“源头考据完成”并置——
 在来源不明时，抑制的对象只能分散到无数二创作者（“法不责众”）；
 来源被指向特定账号后，抑制才有了可执行的目标。
 这一并置目前只有 \conf{C} 级支持，
 本书只把它作为“待检验的时间相关假设”
 保留（第七章 \ref{sec:07:reverse} 的逆反认证机制以此为经验入口）。
\end{enumerate}

“潘多拉奶盒”这个社群用语本身也值得记录：它把“不可逆扩散”
 表达为“盒子打开了就关不上”，与“正式脱离了奶龙”的判断合起来，
 构成了民间叙事中关于“奶蛙—奶龙”关系的 \keyw{断裂叙事}——
 与正典产业线（第三章：授权、维权、关评）构成的“连续叙事”相对照。
 第九章将论证：两种叙事不是对错关系，而是分别描述了“法律层面”
 与“意义层面”的正典—变种关系（第二章定义 2.1 的两个层面）。

\section{考据伦理：关于“创作者”的写作边界}\label{sec:04:ethics}

本章最后一节处理
一条\keyw{硬红线}。

奶蛙起源考据指向了一个具体的、匿名的个体账号。考据完成之后，
 围绕“创作者”的叙述出现了两种危险倾向，本书都必须拒绝：

\begin{enumerate}[label=\arabic*.]
  \item \textbf{“认领叙事”的
伦理风险}：考据指向某个账号后，
 舆论场会自然地要求“当事人出面认领”，
 或在当事人未认领时进行身份推断。本书的立场：
 考据指向的是 \keyw{账号} （一个可公开检索的网络对象），
 不是 \keyw{自然人}。本书不对账号持有者的现实身份作任何推断，
 不呼吁任何个体“认领”或“否认”，
 也不报道任何关于其个人境况（包括健康、职业、情绪、生死）
 的推测性内容——二手评述中确实存在以“最残酷的结局”为标题、
 含个人境况推测的叙述 \conf{C} \cite{shoutanji2026naiwa}，
 本书 \keyw{不转述} 其内容，只在此说明其存在及本书的处理方式。
  \item \textbf{“工具化”的
研究伦理风险}：反过来，也不能因为当事人匿名，就把“创作者”
 当作纯粹的证据节点使用。
 第二章已说明奶蛙的署名权归属原始创作者（附录 D）；
 本书在全部叙述中使用“考据指向的创建者”这一措辞，
 保留其作为 \keyw{权利主体} 的地位，
 而不只作为 \keyw{证据节点} 的地位。
\end{enumerate}

\begin{nlmethod}[伦理操作的三条规则]
\textbf{规则 E1（对象限定）}：考据叙述的宾语只能是账号、智能体、
 帖子、截图等可公开检索对象；不得以“创作者本人”
 为宾语陈述任何个人信息。 \textbf{规则 E2（推测隔离）}：
 对个人境况的推测性材料（无论来自媒体还是社群），
 一律“说明存在、不转述内容”处理。 \textbf{规则 E3（权利保留）}：
 涉及二创形象时，同时写明署名权归属，使“分析”
 不滑向“消费”。
\end{nlmethod}

这三条规则在第十二章会进入“可复现性清单”，
 成为修订者必须遵守的操作条款。在此提前说明，
 是因为本章的全部材料正是它们的 \keyw{第一个适用案例}。

\section{本章小结}\label{sec:04:summary}

本章把奶蛙起源考据拆成了三个阶段：
 误指（ 2024-03 “季洪行说”）、线索（ 2024-12 星野水印）、
 发现与证伪并存（ 2026-08 智能体确认，
 ID \fname{zNCvQDKfXy} 已不可访问）。
 核心结论是一个 \keyw{降级的指向}：
 奶蛙最可能的来源指向星野平台 2024-10-02 的智能体头像，
 该指向由两个 C 级来源交叉印证，但因载体消失而不可复核，
 且考据材料自身声明“仅为推测、缺乏直接证据”。
 音源证据链（王麻衣连麦）作为第二条独立链单独记录，
 级别同为 C 级。 2026-03 的关评拉黑事件是两条产业线的分界，
 本书从中提取了“抑制不彻底”与“指向改变动机”
 两个待检验判断。最后，
 本章给出了三条考据伦理规则（ E1/E2/E3），并以奶蛙“创作者”
 为第一个适用案例——这是全书证据分级制度最严格的一次执行，
 也是第十二章方法论的预演。遗留问题：
 起源考据回答了“形象从哪来”（部分地），
 但没有回答“为什么是这个畸变被选中、
 为什么是这个笑声被绑定”——那是符号问题与机制问题，第六、
 七章的主题。

\begin{nlexercises}
  \item 按定义~\ref{def:04:kaoju} 的结论形态区分，改写以下两句：
 “奶蛙来自星野 AI ”
 与“奶蛙最可能的来源指向星野 AI 的某智能体，
 指向目前不可复核”。说明两句在证据要求上的差异，
 以及为什么第一句在当前材料下不成立。
  \item 给
出一个你熟悉的其他网络梗的“误指 → 线索 → 发现”
 三阶段重构（若该梗起源已公认，
 说明哪一阶段被跳过及原因）。
  \item 规则
E2 要求“说明存在、不转述内容”。
 设计一条可执行的检查程序（如关键词清单 + 句式模板），
 使一位修订者能在 5 分钟内自查修订稿是否违反 E2。
\end{nlexercises}

\nlkeywords 考据、未结案证据链、误指、不可复核性、音源、
关评拉黑、考据伦理、E1/E2/E3
